\scanpage{13}
la deuxième \deleted{\unsure{à|se}} se ramène à
\[
E_n={}_{f^n}E_{n+1}
\]
qui par $\mathrm a_n)$ équivaut à $\supset$, donc à ceci

Soient $x\in E_n$, $h\in A$, \deleted{$f^nx$.} tels qu'il existe $\lambda\in A$, avec
\[
f^n(x,h)=\lambda(\xi_n,f)
\]
alors on a $h\in fA$. En effet, on aura $f^nh=\lambda f$ d'où $\lambda=f^{n-1}h$, (\uline{car $f$ est $A$-régulier}) puis $f^nx=\lambda\xi_n$ \deleted{\unsure{i.e.} $f^{\unsure{n}}$.} i.e.
\[
0=f^{n-1}h\xi_n,\quad\text{or}\quad (\overbrace{f^{n-1}\xi_n}^{y_n})\in {}_fE_n\quad\text{par }\mathrm a_n),
\]
\deleted{\unsure{d'où}} et $ {}_fE_n\simeq A/fA=B$ par $\mathrm b_n)$, \deleted{\unsure{donc}} enfin par construction on a $\xi_n\notin{}_{f^{n-1}}E_n$ i.e. \deleted{$f^{n}$} $y_n\neq0$, donc $hy_n=0\Longrightarrow h\in fA$, cqfd.

Enfin, vérifions $\mathrm c_n$, on a
\[
\begin{aligned}
E_{n+1}/fE_{n+1}&\simeq(E_n\times A)/f(E_n\times A)+A(\xi_n,f)\\
&\simeq(E_n/fE_n\times A/fA)/A(\dot{\xi}_n,\dot f)\\
&\simeq\bigl([B\times(B/gB)^{n-1}]\times B\bigr)/B(g,0,0)\\
&\simeq B/gB\times(B/gB)^{n-1}\times B\simeq B\times(B/gB)^n.
\end{aligned}
\]
% archive sheet number: 12

